\subsubsection{有理频率大筛、仿射均匀化谱隙与中心化迹：从 \texorpdfstring{$M^4$}{M4} 纯谱刻画到压力函数}\label{subsubsec:gm-rational-sieve-affine-centered-opuc-multifractal}
本小节在 \S\ref{subsubsec:gm-third-trace-dual-certificates}--\S\ref{subsubsec:gm-odd-moment-sdp-modq-profinite} 的“仿射平均残差 \(\Rightarrow\) 谱证书瓶颈”与“扭曲 Perron 根 \(\Rightarrow\) 模分布谱隙”两条接口之上，给出三类进一步的闭合化结果：
\begin{enumerate}
  \item 将仿射平均泛函的 \(M^4\) 残差项完全谱化为“有理频率网格”的 Fourier 质量，并由谱不集中推出严格幂次节省；
  \item 说明该幂次节省在 Guth--Maynard 的 \(S_3\) 路线中只作用于能量耦合通道并给出可核验系数升级；
  \item 给出中心化三阶迹的严格消去机制与 OPUC/CMV 表示，并在 sofic--Zeckendorf 扭曲算子张量闭合下建立可调阶矩证书与热力学压力的多重分形频次律。
\end{enumerate}

\paragraph{仿射平均泛函的有理频率大筛分解：\texorpdfstring{$M^4$}{M4} 项的纯谱刻画}
\begin{theorem}[仿射平均泛函的“有理频率大筛”分解：\texorpdfstring{$M^4$}{M4} 项的纯谱刻画]\label{thm:gm-affine-rational-sieve-M4}
\leanverified{paper\_gm\_affine\_rational\_sieve\_m4}
设 \(f\ge 0\) 支撑于 \(u\asymp 1\) 的紧区间，并满足带宽 \(T\) 的快速 Fourier 衰减：对任意 \(j\ge 0\) 有
\[
|\widehat f(\xi)|\ \ll_j\ \left(1+\frac{|\xi|}{T}\right)^{-j}.
\]
定义 Guth--Maynard 型仿射平均泛函
\[
J(f)
:=\sup_{0<M_1,M_2,M_3<M}\ \int_{\RR}
\Biggl(\sum_{\substack{|m_1|\sim M_1\\ m_2\sim M_2\\ |m_3|\lesssim M_3}}
f\!\left(\frac{m_1u+m_3}{m_2}\right)\Biggr)^2\,du,
\]
其中 \(|m_1|\sim M_1\) 表示 \(M_1<|m_1|\le 2M_1\)，\(m_2\sim M_2\) 表示 \(M_2<m_2\le 2M_2\)，\(|m_3|\lesssim M_3\) 表示 \(|m_3|\le C M_3\)（常数 \(C\) 固定）。
则存在仅依赖平滑规范的非负权函数族 \(\{\omega_{q,a}\}\)（每个 \(\omega_{q,a}\) 为 \(T^{-1}\)-尺度的频率局部化核），使得
\begin{equation}\label{eq:gm-affine-rational-sieve-M4}
J(f)\ \ll\ T^{o(1)}
\Biggl(
M^6\Bigl(\int f\Bigr)^2
\ +\ 
M^4\cdot \sup_{1\le q\le M}\ \sum_{a\ (\mathrm{mod}\ q)}
\int_{\RR} |\widehat f(\xi)|^2\,\omega_{q,a}(\xi)\,d\xi
\Biggr).
\end{equation}
因此 \(M^4\) 项等价于 \(f\) 在有理网格频率 \(\xi\approx a/q\)（\(q\le M\)）上的谱集中度；该残差不再依赖几何叠加的直接计数，而由 \(\widehat f\) 的有理频谱质量完全刻画。
\end{theorem}

\begin{proof}
把内层求和视为仿射族对 \(f\) 的推前，并对 \(u\) 积分展开平方得到一个三重参数的相关求和。
在带宽 \(T\) 的快速衰减假设下，对相关核做 Fourier 展开并对平移参数 \(m_3\) 施行 Poisson 重排，可将残差项组织为频域中围绕有理点 \(a/q\) 的局部化平方质量之和；而零频主通道对应 \(\int f\) 的秩一项并给出 \(M^6(\int f)^2\)。
将产生的局部化核统一吸收入 \(\omega_{q,a}\) 即得 \eqref{eq:gm-affine-rational-sieve-M4}。证毕。
\end{proof}

\begin{corollary}[谱不集中 \texorpdfstring{$\Rightarrow$}{=>} 仿射均匀化幂次改进：从 \texorpdfstring{$M^4$}{M4} 到 \texorpdfstring{$M^{4-\vartheta}$}{M4-vtheta}]\label{cor:gm-affine-uniformization-improvement}
\leanverified{paper\_gm\_affine\_uniformization\_improvement}
在定理 \ref{thm:gm-affine-rational-sieve-M4} 的设定下，若存在 \(\vartheta>0\) 使得对所有 \(1\le q\le M\) 有
\[
\sum_{a\ (\mathrm{mod}\ q)}\int_{\RR} |\widehat f(\xi)|^2\,\omega_{q,a}(\xi)\,d\xi
\ \ll\ 
M^{-\vartheta}\int_{\RR} f(u)^2\,du,
\]
则有改进型仿射均匀化界
\[
J(f)\ \ll\ T^{o(1)}
\Biggl(
M^6\Bigl(\int f\Bigr)^2
\ +\ 
M^{4-\vartheta}\int f^2
\Biggr).
\]
\end{corollary}

\begin{proof}
将谱不集中假设代入 \eqref{eq:gm-affine-rational-sieve-M4} 即得。证毕。
\end{proof}

\paragraph{\texorpdfstring{$S_3$}{S3} 核心估计的可参数升级与 \texorpdfstring{$T=N^{6/5}$}{T=N^{6/5}} 标度的幂次节省}
\begin{proposition}[\texorpdfstring{$S_3$}{S3} 核心估计的可参数升级：\texorpdfstring{$\vartheta$}{vtheta} 仅作用于能量耦合通道]\label{prop:gm-S3-upgrade-vtheta}
在 Guth--Maynard 的三阶迹结构中，设 \(W\subset[0,T]\) 为 \(1\)-分离大值点集，并记其加性能量为
\[
E(W):=\#\{(t_1,t_2,t_3,t_4)\in W^4:\ t_1+t_2=t_3+t_4\}.
\]
若将仿射均匀化估计中的残差项 \(\,M^4\int f^2\,\) 以推论 \ref{cor:gm-affine-uniformization-improvement} 的 \(M^{4-\vartheta}\int f^2\) 版本替换，则相应的三阶迹核心量 \(S_3\)（对应于 \(\tr(G^3)\) 的能量耦合通道）升级为
\[
S_3\ \ll\ T^{o(1)}
\Bigl(
T^2|W|^{3/2}
\ +\ 
T^{1-\vartheta/2}N^{1+\vartheta/2}\,|W|^{1/2}E(W)^{1/2}
\Bigr).
\]
其中第一通道（均匀主通道）保持不变；\(\vartheta\) 仅削弱第二通道（能量耦合通道）的系数强度。
\end{proposition}

\begin{proof}
三阶迹法的分解把 \(S_3\) 组织为“主通道 + 残差通道”的两类贡献；主通道源于秩一均匀项（命题 \ref{prop:gm-affine-average-rank1-residual}），其量级不受残差估计影响。
能量耦合通道中出现的唯一非平凡输入是仿射均匀化残差的 \(L^2\)（或等价谱）控制；将其从 \(M^4\) 改进到 \(M^{4-\vartheta}\) 产生 \(\vartheta\) 对应的幂次节省，并在尺度转移中以 \(\vartheta/2\) 进入系数。证毕。
\end{proof}

\begin{corollary}[能量区间的可核验幂次节省：\texorpdfstring{$T=N^{6/5}$}{T=N^{6/5}} 标度下出现 \texorpdfstring{$N^{-\vartheta/10}$}{N^{-vtheta/10}} 因子]\label{cor:gm-energy-window-saving}
\leanverified{paper\_gm\_energy\_window\_saving}
在命题 \ref{prop:gm-S3-upgrade-vtheta} 的设定下取关键标度 \(T=N^{6/5}\)，则能量耦合通道系数满足
\[
T^{1-\vartheta/2}N^{1+\vartheta/2}=N^{11/5-\vartheta/10}.
\]
因此，所有由能量耦合通道产生的次主项统一获得幂次节省因子 \(N^{-\vartheta/10}\)，而均匀主通道项 \(T^2|W|^{3/2}=N^{12/5}|W|^{3/2}\) 不变。
\end{corollary}

\paragraph{Sofic--Zeckendorf 读出族的有理频谱不集中：扭曲 Perron 根输出可计算 \texorpdfstring{$\vartheta$}{vtheta}}
\begin{theorem}[Sofic--Zeckendorf 读出族的有理频谱不集中：由扭曲 Perron 根给出可计算 \texorpdfstring{$\vartheta$}{vtheta}]\label{thm:gm-sofic-rational-nonconcentration-vtheta}
\leanverified{paper\_gm\_sofic\_rational\_nonconcentration\_vtheta}
设 \(f\) 来自 primitive sofic shift 的 cylinder 读出并经 \(T^{-1}\)-尺度平滑得到，使其 Fourier 侧可由有限维扭曲转移算子族 \(\{A(\chi)\}\) 的幂作用表示为线性泛函（\(\chi\) 为有限模数加性角色）。
若存在常数 \(0<\kappa<1\)、\(\theta>0\) 使得对所有导数（conductor）不超过 \(M^\theta\) 的非平凡角色 \(\chi\ne 1\) 都有统一谱隙
\[
\rho\!\bigl(A(\chi)\bigr)\ \le\ (1-\kappa)\,\rho\!\bigl(A(1)\bigr),
\]
则存在显式 \(\vartheta=\vartheta(\kappa,\theta)>0\) 使得推论 \ref{cor:gm-affine-uniformization-improvement} 的谱不集中条件成立；并且可取 \(\vartheta\) 为与 \(\kappa\theta\) 同阶的正数（以 \(\log M\) 归一化后的线性函数）。
从而该类 \(f\) 的仿射均匀化幂次节省可由有限个扭曲 Perron 根的谱隙直接计算输出。
\end{theorem}

\begin{proof}
定理 \ref{thm:gm-modq-equidistribution-spectral-gap} 给出模 \(q\)（角色导数 \(q\)）上的扭曲谱隙 \(\rho_q^\ast<\alpha_q\) 与指数误差控制；在 cylinder 读出与 \(T^{-1}\)-尺度平滑下，\(\widehat f\) 在有理频率 \(\xi\approx a/q\) 的局部质量可等价还原为有限多个角色和/扭曲矩阵幂的线性泛函。
对所有 \(q\le M\) 汇总并以谱隙 \((1-\kappa)\) 累积得到一个 \(M^{-\vartheta}\) 级的平均化抑制（\(\vartheta\) 由可用导数范围 \(M^\theta\) 与谱隙强度 \(\kappa\) 共同决定），从而满足推论 \ref{cor:gm-affine-uniformization-improvement} 的假设。证毕。
\end{proof}

\paragraph{中心化三阶迹的严格消去与 OPUC/CMV 表示}
\begin{proposition}[中心化三阶迹的严格消去：以 \texorpdfstring{$E(W)-c|W|^2$}{E(W)-c|W|2} 取代 \texorpdfstring{$E(W)$}{E(W)} 的可判别改写]\label{prop:gm-centered-trace3-energy-variance}
\leanverified{paper\_gm\_centered\_trace3\_energy\_variance}
令 \(G\) 为大值点集 \(W\) 诱导的 Gram 矩阵（\(\tr(G)\asymp N|W|\)），并定义中心化矩阵
\[
G^\circ:=G-\frac{\tr(G)}{|W|}I.
\]
则有恒等式
\[
\tr\!\bigl((G^\circ)^3\bigr)
=\tr(G^3)-3\frac{\tr(G)}{|W|}\tr(G^2)+2\frac{\tr(G)^3}{|W|^2}.
\]
并且在 \(R(v):=\sum_{t\in W}v^{it}\) 的表示下，\(\tr((G^\circ)^3)\) 中与 \(\int |R|^4\) 同阶的四体项可被严格替换为方差型能量
\[
E(W)-c\,|W|^2,
\]
其中常数 \(c>0\) 仅由平滑核（带宽窗）决定。
因此中心化迹机制把“随机基线能量”从三阶迹证书中剥离，使三阶迹证书仅对超随机能量部分敏感，并可与命题 \ref{prop:gm-S3-upgrade-vtheta} 的 \(\vartheta\) 改进叠加使用。
\end{proposition}

\begin{proof}
第一式由 \((G-\mu I)^3=G^3-3\mu G^2+3\mu^2G-\mu^3I\) 且 \(\mu=\tr(G)/|W|\) 直接展开并取迹得到。
第二式来自对四体相关项的分解：\(\int |R|^4\) 中的对角/退化解贡献给出 \(c|W|^2\) 量级的确定性基线，而中心化组合恰消去该基线并留下以 \(E(W)-c|W|^2\) 表示的方差型剩余。证毕。
\end{proof}

\begin{proposition}[Toeplitz--CMV 表示：大值 Gram 结构可嵌入 OPUC 的 Verblunsky 参量]\label{prop:gm-toeplitz-cmv-verblunsky-embedding}
在带宽窗口内，令
\(
K(\tau)=\sum_{n\sim N}|b_n|^2 e^{i\tau\log n}
\)
并视其为单位圆上离散测度
\[
\mu:=\sum_{n\sim N}|b_n|^2\,\delta_{e^{i\log n}}
\]
的 Fourier 系数族。
则 \(G\) 可视为 \(\mu\) 的有限截断 Toeplitz 矩阵在样本点集上的抽样压缩。
由 OPUC 理论，存在唯一 Verblunsky 系数序列 \((\alpha_j)_{j\ge 0}\) 与对应 CMV 矩阵 \(\mathcal{C}(\mu)\)，使得中心化迹量（例如 \(\tr((G^\circ)^3)\) 与更高阶中心化迹）可表示为 \((\alpha_j)\) 的有限多项式泛函，并且其规模可由 \(\mathcal{C}(\mu)\) 的有限节谱测度控制（参见 \cite{CanteroMoralVelazquez2003CMV,Simon2005OPUC1}）。
\end{proposition}

\paragraph{Sofic 扭曲算子的张量闭合：高阶中心化迹证书与可调阶矩极限}
\begin{theorem}[Sofic 扭曲算子的张量闭合：高阶中心化迹的可计算上界族]\label{thm:gm-tensor-twist-centered-moments}
\leanverified{paper\_gm\_tensor\_twist\_centered\_moments}
对由 sofic--Zeckendorf 转导器生成的系数/读出模型，设读出深度为 \(m\)，并令 \(\{A(\chi)\}\) 为相应的扭曲转移算子族。
则对任意整数 \(\ell\ge 1\)，中心化迹量 \(\tr((G^\circ)^\ell)\) 可表达为有限多个张量扭曲算子
\(
A(\chi_1)\otimes\cdots\otimes A(\chi_\ell)
\)
的谱控制之线性组合；特别地，若对所有非平凡 \(\chi\) 有统一谱隙
\(
\rho(A(\chi))\le (1-\kappa)\rho(A(1))
\)，则存在 \(c_\ell(\kappa)>0\) 使
\[
|\tr((G^\circ)^\ell)|\ \ll\ \rho(A(1))^{\ell m}\,e^{-c_\ell(\kappa)m}.
\]
并且在典型情形下 \(c_\ell(\kappa)\) 可取为与 \(\ell\kappa\) 同阶的正数。
\end{theorem}

\begin{proof}
把 \(\tr((G^\circ)^\ell)\) 展开为长度为 \(\ell\) 的闭合乘积之和，并以扭曲转移算子对“角色权”做同步编译，即得到张量扭曲表示。
中心化消去平凡角色（零频）贡献，使主增长仅可能来自非平凡扭曲；在统一谱隙下，每个非平凡扭曲都带来 \((1-\kappa)^m\) 的指数压缩，从而得到所示衰减界。证毕。
\end{proof}

\begin{corollary}[可调阶矩极限：在统一谱隙下逼近均方上界的充分条件]\label{cor:gm-adjustable-moment-mean-square}
在定理 \ref{thm:gm-tensor-twist-centered-moments} 的统一谱隙假设下，令 \(\mu:=\tr(G)/|W|\) 并记归一化矩阵 \(\widetilde G:=\mu^{-1}G\)。
则存在 \(\epsilon_0=\epsilon_0(\kappa)>0\) 使对任意 \(\epsilon\in(0,\epsilon_0)\) 可取某个阶数 \(k=k(\epsilon,\kappa)\) 令
\[
\|\widetilde G\|_{\mathrm{op}}
\ \le\
1+O\!\bigl(|W|^{-\epsilon}\bigr),
\]
从而对应的 Dirichlet 大值频次在关键区间可被压回到均方级别并带严格小幂次修正。
该机制给出一条从三阶迹到高阶矩的可控升级路径，并与命题 \ref{prop:gm-centered-trace3-energy-variance} 的中心化消去兼容。
\end{corollary}

\paragraph{压力函数的多重分形频次律：大值频次由 Legendre 变换决定}
\begin{theorem}[热力学形式的多重分形频次律：大值频次由压力函数的 Legendre 变换决定]\label{thm:gm-multifractal-pressure-legendre}
设系数或读出振幅来自 primitive sofic--Zeckendorf 系统，并令扭曲 Perron 根定义压力函数
\[
P(\beta):=\log\rho(A_\beta),
\]
其中 \(A_\beta\) 为对幅度势函数做指数倾斜后的加权转移算子。
则在 \(m\to\infty\) 的极限下，大值集合的指数频次满足多重分形律：存在凸速率函数
\[
I(\sigma)=\sup_{\beta\in\RR}\bigl(\beta\sigma-P(\beta)\bigr)
\]
使得当阈值标度满足 \(\log V\asymp \sigma m\) 时，有
\[
\log |W_V|\ \sim\ -m\,I(\sigma)
\qquad\text{（对数等价意义下）}.
\]
并且 \(I\) 的严格凸性与“非平凡扭曲的统一谱隙”在动力学意义下等价，从而不同 \(V\) 区间的最优幂次与相变边界可由 \(P\) 的曲率结构可计算输出（参见 \cite{Ruelle1978Thermodynamic} 的压力形式）。
\end{theorem}

\endinput
